\section{Results}\label{sec:results}

\subsection{Heat damages cereal yields}

Table~\ref{tab:firststage} shows that the hottest-month temperature anomaly lowers cereal-yield
growth substantially. With country and year fixed effects, a 1$^{\circ}$C hotter-than-trend
hottest month reduces yield growth by 3.6 percentage points across 155 countries (SE 0.53), and by
2.7 points in emerging and developing economies (SE 0.76). The annual mean anomaly has no
detectable average effect ($-0.61$, SE 1.01), in line with the threshold effects of heat documented
by \citet{schlenker2009}. We report all three shock definitions we examined, because the choice of
the hottest-month anomaly was made after seeing this table. Among countries whose baseline climate
is already hot ($\geq 20^{\circ}$C), the estimate is imprecise ($-0.34$, SE 0.91), so the
first stage is driven mainly by cooler and temperate producers. This matters for interpreting the
allocation rules below.

\begin{table}[t]
\centering
\caption{Heat shocks and cereal-yield growth: first stage}\label{tab:firststage}
\small
\begin{tabular}{lcccc}
\toprule
Sample & Coefficient (SE) & $p$ & Observations & Countries \\
\midrule
All countries & $-3.572$ (0.530) & $<0.001$ & 3,679 & 155 \\
EMDEs & $-2.684$ (0.759) & $<0.001$ & 2,538 & 107 \\
Baseline temp.\ $\geq 20^{\circ}$C & $-0.336$ (0.908) & 0.711 & 1,971 & 83 \\
Annual-mean anomaly as regressor & $-0.612$ (1.006) & 0.543 & 3,679 & 155 \\
\bottomrule
\end{tabular}
\begin{minipage}{0.95\textwidth}\footnotesize\smallskip
\emph{Notes:} Dependent variable: $100\times\Delta\ln$ cereal yield (kg/ha). Regressor: hottest-month
mean temperature minus its 1980--2023 country-specific linear trend ($^{\circ}$C), except in the last
row. Country and year fixed effects; standard errors clustered by country. Observations with
$|\Delta\ln\text{yield}|\geq 1$ are dropped as data errors. The sample is larger than the estimation
sample of Table~\ref{tab:ate}, which also requires the lagged covariates.
\end{minipage}
\end{table}

\subsection{No detectable resilience dividend}

Table~\ref{tab:ate} reports the main estimates for 3,336 country-years in 145 countries. The
median across the 20 sample splits of the average resilience effect is 1.02 percentage points of
yield growth per $^{\circ}$C, with a split-adjusted standard error of 2.12. The 95 per cent
confidence interval, $[-3.13, 5.17]$, is wide. It contains zero, and it also contains effects
larger than the average damage itself. The data therefore cannot distinguish between climate-AI
capability having no effect on heat damage and its offsetting all of it. Other ways of computing
the estimate give the same picture. Using the observation-level median score, the estimate is 0.84
with a country-clustered SE of 1.01, a two-way (country and year) clustered SE of 0.94, or a
Driscoll--Kraay SE of 1.17. Winsorising the scores at the 1st and 99th percentiles gives 0.51
(SE 0.74). The scores are heavy-tailed: the 99th percentile of $|\Gamma|$ is about 300 against a
median of 11.

The average marginal damage implied by the orthogonal scores is similar in the two groups
($-1.49$ with capability and $-1.31$ without, split-adjusted; $-1.52$ overall). That is less than half the fixed-effects estimate on the
same sample ($-3.58$ with country and year effects, and $-3.17$ with year and region effects and the
same covariates; Appendix Table~\ref{tab:recon}). Country fixed effects are therefore not the
reason. The difference arises within the orthogonal estimator, whose scores weight observations by
the inverse of the estimated conditional variance of the shock. We have not been able to reconcile
the two fully, and we treat the scale of the orthogonal estimates with caution.

The classical specification again behaves differently. A two-way fixed-effects regression with a
shock-by-capability interaction gives 0.32 (SE 0.97) for yields. For agricultural value-added growth
it gives 0.82 (SE 0.46, $p\approx0.07$), which would often be reported as weak evidence of a
resilience dividend. The doubly robust estimate for the same outcome is 0.07 (split-adjusted SE
1.01). Once the flexible relationship between capability, income, connectivity and climate is taken
into account, the suggestive interaction disappears, as the Monte Carlo evidence on the interaction
model would lead one to expect.

\begin{table}[t]
\centering
\caption{Resilience effect of climate-AI capability on cereal-yield growth}\label{tab:ate}
\small
\begin{tabular}{lcc}
\toprule
 & Estimate (SE) & 95\% CI \\
\midrule
\multicolumn{3}{l}{\emph{EW-DRPL, average resilience effect}}\\
Median over splits, split-adjusted SE (main) & 1.022 (2.118) & [$-3.13$, 5.17] \\
Observation-level median score, country clusters & 0.843 (1.014) & [$-1.14$, 2.83] \\
\quad two-way clusters (country, year) & 0.843 (0.937) & [$-1.00$, 2.68] \\
\quad Driscoll--Kraay (2 lags) & 0.843 (1.172) & [$-1.46$, 3.14] \\
\quad scores winsorised at 1st/99th percentiles & 0.505 (0.738) & [$-0.94$, 1.95] \\
Split-to-split range & [$-2.56$, 7.47] & \\
\midrule
\multicolumn{3}{l}{\emph{EW-DRPL, mean marginal damage}}\\
Capability ($T=1$, 1,734 obs.) & $-1.489$ (1.188) & [$-3.82$, 0.84] \\
No capability ($T=0$, 1,602 obs.) & $-1.310$ (1.214) & [$-3.69$, 1.07] \\
\midrule
\multicolumn{3}{l}{\emph{Two-way fixed effects (for comparison)}}\\
Heat shock & $-3.386$ (0.820) & \\
Heat shock $\times$ capability & 0.315 (0.970) & \\
Heat shock $\times$ capability, agricultural value added & 0.817 (0.457) & \\
\midrule
\multicolumn{3}{l}{\emph{EW-DRPL, agricultural value added}}\\
Average resilience effect (split-adjusted) & 0.066 (1.014) & [$-1.92$, 2.05] \\
\bottomrule
\end{tabular}
\begin{minipage}{0.95\textwidth}\footnotesize\smallskip
\emph{Notes:} Units: percentage points of annual growth per $^{\circ}$C of hottest-month anomaly.
3,336 country-years in 145 countries (3,305 for agricultural value added). Capability:
climate-AI knowledge stock above the level predicted by income in the previous year. Covariates,
lagged one year: log GDP per capita, internet use, agricultural share, urban share, log population,
log fertiliser use, ND-GAIN readiness and general AI output; also baseline temperature, the
precipitation shock, year and region. Damage rows: split-adjusted SEs.
\end{minipage}
\end{table}

\subsection{Heterogeneity}

Neither the best linear projection of $\Gamma$ on the policy features nor the sorted group effects
reveal a reliable gradient. All six projection coefficients are within 1.4 standard errors of zero,
the largest relative to its SE being inequality ($-1.63$, SE 1.23) and knowledge sovereignty
($-1.58$, SE 1.47). The difference between the top and bottom quintiles of predicted resilience is
4.5 (SE 3.7), and the quintile means are not monotone (Appendix Table~\ref{tab:gates}). The
classification analysis shows that the model predicts the largest resilience effects for
country-years with lower internet use (22 percentage points lower in the top than in the bottom
quintile, SE 4.4), lower knowledge sovereignty and lower income. Because the sorted effects are not
monotone, this describes where the prediction model places resilience rather than where resilience
is. With about 145 countries and noisy scores, the data cannot rank countries by the size of the
resilience dividend.

\subsection{Allocation rules}

Table~\ref{tab:policy} reports the held-out welfare values of the allocation rules for a budget of
30 per cent of the 126 countries with complete distributional data. Values are means over countries
of $\pi(Z_i)B_i(\epsilon)$, which combine the estimated resilience scores with shock size,
population, agricultural exposure and welfare weights.

The table supports three conclusions, and none of them favours data-driven targeting in this
application. First, no rule's value is distinguishable from zero, which is what one would expect
given that the average resilience effect is not. Second, simple rules do at least as well as the
learned trees. Allocating capacity to the poorest 30 per cent of countries has the highest point
value under every $\epsilon>0$, and at $\epsilon=2$ it beats the EW-DRPL tree by 1.39
[$0.02$, $3.26$], the only policy comparison whose interval excludes zero. Third, the rule closest
to building on existing strength, which allocates capacity to the countries with the largest
climate-AI stock, has a value of essentially zero. This is largely mechanical: those countries have
little agricultural exposure and comparatively rich deciles, so even a large resilience effect would
add little equity-weighted welfare. Its difference from the poorest-first rule is not
statistically significant ($-0.65$ [$-1.80$, $0.31$] at $\epsilon=1$).

The equity gain of the $\epsilon=1$ rule over the utilitarian rule is 0.33 [$-0.01$, 0.90] and the
equity price is $-0.06$ [$-0.16$, 0.01]. A negative price, meaning that the equity-weighted rule also
does better on utilitarian welfare, cannot be a property of the population optimum. It reflects the
noise of cross-fitted trees. When the objective counts countries rather than people, every rule's
value is close to zero, and the advantage of poorest-first disappears (Appendix
Table~\ref{tab:unweighted}). The population-weighted results are thus driven by a few populous,
poor and agricultural countries.

The learned trees are readable, but they are not stable. For $\epsilon\le1$ the tree treats 31
countries. Nineteen are low-income countries with a hot baseline climate: eighteen in sub-Saharan Africa
and Myanmar. Twelve have higher incomes and a predominantly domestic climate-AI research base,
including China, India, Indonesia, the Philippines and the United States. At $\epsilon=2$ the root
split moves to internet use, and the low-income branch shrinks to ten countries. Viet Nam is not
treated by any rule; its estimated equity-weighted benefit is slightly negative at $\epsilon=1$
and $\epsilon=2$.
The hot, low-income branch is also the group for which the first stage is weakest
(Table~\ref{tab:firststage}), so its selection reflects exposure and welfare weights more than any
measured resilience effect.

\begin{table}[t]
\centering
\caption{Held-out welfare value of allocation rules (budget $k=0.3$)}\label{tab:policy}
\small
\begin{tabular}{lcccc}
\toprule
Rule & $\epsilon=0$ & $\epsilon=0.5$ & $\epsilon=1$ & $\epsilon=2$ \\
\midrule
EW-DRPL ($\epsilon$-specific tree) & 0.190 & 0.344 & 0.399 & $-0.119$ \\
 & {\footnotesize[$-0.17$, 0.74]} & {\footnotesize[$-0.25$, 1.13]} & {\footnotesize[$-0.40$, 1.30]} & {\footnotesize[$-1.10$, 0.84]} \\
Utilitarian tree & 0.190 & 0.185 & 0.072 & $-0.282$ \\
 & {\footnotesize[$-0.17$, 0.74]} & {\footnotesize[$-0.37$, 0.98]} & {\footnotesize[$-0.69$, 1.01]} & {\footnotesize[$-1.36$, 0.72]} \\
Plug-in tree & 0.387 & 0.284 & 0.475 & 0.789 \\
 & {\footnotesize[$-0.05$, 0.98]} & {\footnotesize[$-0.22$, 1.06]} & {\footnotesize[$-0.19$, 1.37]} & {\footnotesize[$-0.22$, 2.08]} \\
Capable-first & $-0.006$ & $-0.007$ & $-0.007$ & $-0.009$ \\
 & {\footnotesize[$-0.02$, 0.01]} & {\footnotesize[$-0.02$, 0.01]} & {\footnotesize[$-0.04$, 0.02]} & {\footnotesize[$-0.05$, 0.02]} \\
Poorest-first & 0.310 & 0.471 & 0.647 & 1.268 \\
 & {\footnotesize[$-0.08$, 0.89]} & {\footnotesize[$-0.17$, 1.35]} & {\footnotesize[$-0.33$, 1.79]} & {\footnotesize[$-0.62$, 3.26]} \\
Most-agricultural-first & 0.291 & 0.431 & 0.581 & 1.169 \\
 & {\footnotesize[$-0.11$, 0.87]} & {\footnotesize[$-0.21$, 1.30]} & {\footnotesize[$-0.41$, 1.70]} & {\footnotesize[$-0.64$, 3.16]} \\
Random (expected) & 0.130 & 0.176 & 0.213 & 0.358 \\
\midrule
Equity gain $W_\epsilon(\hat\pi_\epsilon)-W_\epsilon(\hat\pi_0)$ & & & 0.327 & 0.164 \\
 & & & {\footnotesize[$-0.01$, 0.90]} & {\footnotesize[$-0.04$, 0.52]} \\
Equity price $W_0(\hat\pi_0)-W_0(\hat\pi_\epsilon)$ & & & $-0.061$ & $-0.020$ \\
 & & & {\footnotesize[$-0.16$, 0.01]} & {\footnotesize[$-0.09$, 0.03]} \\
EW-DRPL minus poorest-first & & & $-0.248$ & $-1.386$ \\
 & & & {\footnotesize[$-0.77$, 0.14]} & {\footnotesize[$-3.26$, $-0.02$]} \\
\bottomrule
\end{tabular}
\begin{minipage}{0.95\textwidth}\footnotesize\smallskip
\emph{Notes:} 126 countries, of which 86 EMDEs; capacity for 37. Each rule is learned on four
country folds and evaluated on the fifth; columns give the inequality aversion used in evaluation.
Values: mean over countries of $\pi(Z_i)B_i(\epsilon)$, where $B_i$ includes the shock size, the
population scale and the agricultural exposure of each decile (allocated to the lowest deciles
first). 95 per cent country-cluster bootstrap intervals in brackets (500 replications), holding the
learned rules fixed. Capable-first: largest climate-AI stock per capita; poorest-first: lowest GDP
per capita; most-agricultural-first: highest agricultural employment share. Trees treat fewer than
37 countries because a depth-2 tree can only treat whole leaves.
\end{minipage}
\end{table}

\subsection{Robustness}\label{sec:robust}

Table~\ref{tab:sens} reports alternative specifications. The average resilience effect stays
imprecise and close to zero relative to its standard error when the non-orthogonal score is used,
when propensity scores are clipped at $[0.05,0.95]$, with random-forest nuisance models, in the EMDE
subsample, with the capability indicator that is not adjusted for income, and for agricultural value
added. The exception is the annual-mean anomaly, for which the estimate is $-5.59$ (SE 1.92). The
annual anomaly has no first-stage effect on yields on average (Table~\ref{tab:firststage}), so this
estimate compares sensitivities that are both close to zero. It is nevertheless the one
specification in which capability is associated with significantly greater sensitivity, and we
report it as such.

The allocation rules are much less stable than the average effect. Relative to the main rules, the
treated sets are unchanged under alternative clipping, and at $\epsilon=2$ also under uniform
exposure. Under the other specifications they overlap only partly, with Jaccard indices between
0.02 and 0.71 at $\epsilon=1$ and between 0.08 and 0.51 at $\epsilon=2$. With a budget of 20 per
cent, the held-out values of the learned rules lie between $-0.28$ and 0.08. This instability is the empirical
counterpart of the small-sample regret documented in the Monte Carlo study, and it is why we do not
draw country-specific recommendations from the trees.

\begin{table}[t]
\centering
\caption{Sensitivity analysis}\label{tab:sens}
\small
\resizebox{\textwidth}{!}{%
\begin{tabular}{lccccc}
\toprule
 & & & \multicolumn{3}{c}{Treated set vs main} \\
\cmidrule(lr){4-6}
Specification & Obs. & Average effect (SE) & $n$ ($\epsilon=1$) & Jaccard ($\epsilon=1$) & Jaccard ($\epsilon=2$) \\
\midrule
Main & 3,336 & 0.843 (1.014) & 31 & 1.00 & 1.00 \\
(a) Non-orthogonal slope score & 3,336 & 0.727 (1.126) & 22 & 0.71 & 0.51 \\
(b) Clipping at $[0.05,0.95]$ & 3,336 & 0.978 (1.070) & 31 & 1.00 & 1.00 \\
(c) Annual-mean anomaly as shock & 3,336 & $-5.594$ (1.916) & 33 & 0.21 & 0.11 \\
(d) Random-forest nuisance models & 3,336 & $-0.138$ (0.760) & 22 & 0.71 & 0.42 \\
(e) EMDEs only & 2,247 & 0.669 (1.417) & 25 & 0.51 & 0.28 \\
(f) Capability without income adjustment & 3,336 & 2.258 (1.672) & 33 & 0.14 & 0.09 \\
(g) Agricultural value added as outcome & 3,305 & 0.001 (0.610) & 26 & 0.02 & 0.08 \\
(h) Uniform exposure across deciles & -- & -- & 22 & 0.71 & 1.00 \\
\bottomrule
\end{tabular}}
\begin{minipage}{0.95\textwidth}\footnotesize\smallskip
\emph{Notes:} Average effect: observation-level median score over 20 splits, country-clustered SE.
$n$: number of countries treated by the full-sample $\epsilon=1$ tree. Jaccard: overlap with the
main treated set.
\end{minipage}
\end{table}

\subsection{Decile incomes, aggregate growth and the cost of AI}\label{sec:res_decile}

We also applied EW-DRPL directly to decile income growth from 1,334 household-survey spells in 124
countries, and to GDP per capita growth (Appendix~\ref{app:decile}). Neither has a usable first
stage. The annual temperature anomaly has no detectable effect on GDP growth in our 1998--2023
sample ($-0.10$, SE 0.18), even in hot EMDEs ($-0.27$, SE 0.46), and spell-averaged anomalies have
none on decile growth. Without measurable damage there is nothing for capability to reduce, and the
resulting resilience estimates are uninformative.

On the cost side, double machine learning estimates of the association between the climate-AI
knowledge stock and per-capita electricity use ($-0.015$, SE 0.112) and CO$_2$ emissions (0.061,
SE 0.099), and between general AI output and the same outcomes (0.070, SE 0.060, and 0.028, SE
0.056), are imprecise and centred near zero. These intervals do not rule out economically meaningful
effects, but they give no indication that national AI research capacity drives economy-wide energy
demand. The electricity footprint documented by \citet{iea2025} arises where data centres are
located, which is not necessarily where AI research is produced.
